\section{来源审计：这不是一堂公司宣传课}

这段访谈表面上按时间讲述 Guillaume Lample 从无监督机器翻译、形式化数学、LLaMA 到 Mistral 的经历；真正值得学习的却不是履历，而是\textbf{研究问题如何不断改变系统边界}。当任务从“对齐两个表示空间”变成“搜索可验证证明”，再变成“训练并部署通用模型”，团队所需要的数据、反馈信号、计算基础设施和产品接口都会改变。因此，本讲的核心问题是：\textbf{一个模型何时只是研究产物，何时才成为可运行、可定制、可治理的系统？}

历史公开视频在 2026 年 8 月 11 日核验时已经转为 private；本地仍保存官方封面和完整时间戳字幕。字幕是自动识别稿，包含人名和产品名错误，因此讲义只在上下文明确时做规范化，并以论文或官方页面校验算法、模型和监管事实。课堂中的人员规模、成本、下载量等数字统一视为\textbf{讲者在当时的口述估计}，而不是静态公司事实。

\begin{importantbox}{整讲主线：四次系统边界扩张}
\begin{enumerate}
  \item 无监督翻译把问题定义为\textbf{表示空间与学习闭环}；
  \item 形式化数学把问题定义为\textbf{环境、动作、验证器与搜索}；
  \item LLaMA 把问题定义为\textbf{数据规模、数值稳定性与推理成本}；
  \item Mistral 把问题定义为\textbf{部署、定制、反馈、隐私与持续运营}。
\end{enumerate}
\end{importantbox}

带着这条主线读图时，不要把每个阶段看成孤立项目。它们之间的箭头表示“前一阶段暴露的新瓶颈”，而不是简单的职业跳槽时间线。

\lecturefigure{01-research-trajectory.png}{研究问题逐步变成新的基础设施要求。}{本地时间戳字幕 03:00--18:00；\texttt{lecture04-diagrams.py} 可复现重绘。}

\begin{knowledgebox}{读图：先找每一阶段的权威反馈信号}
无监督翻译没有人工平行语料，只能依靠表示结构和循环重建；形式化证明的权威来自 proof assistant，而不是模型自评；预训练依靠 held-out loss 与下游评测；产品系统还必须接受延迟、失败率、用户反馈、数据边界和成本约束。\textbf{反馈信号越接近真实结果，系统就越能进行可靠迭代。}
\end{knowledgebox}

\begin{warningbox}{open source、open weights 与可用系统不是同义词}
社区经常把可下载权重统称为“开源模型”。本讲更精确地区分：\term{open weights} 是可获得参数；\term{open license} 决定可修改和商用的边界；\term{reference code} 提供基础运行实现；\term{production system} 还需要 serving、工具、评测、安全、监控和更新机制。拿到 checkpoint 并不等于拿到产品。
\end{warningbox}

\subsection{本章小结}

本讲用“系统边界扩张”而不是“公司发展史”组织材料。后续每一节都要同时回答三个问题：任务状态是什么，正确性由谁判断，模型之外还缺哪些系统部件。

